\scanpage{438}
\mspage{5}
% archive sheet number: 437

Pour la flèche $f_! f^{\unsure{\prime|*}}H\longrightarrow H$ sur $T$, il faut trouver pour \unsure{tt} $t\in\mathrm{Ob}\,T$, et fonctoriellement en $t$

\diagram{\begin{tikzcd}[ampersand replacement=\&,row sep=4mm]
\underset{\raisebox{0pt}{\rule{.35pt}{.5ex}\rule{3.5em}{.35pt}\rule{.35pt}{.5ex}}}{(f_!f^*H)_t}\arrow[r]\arrow[d,equals,"\mathrm{def}"] \& H_t\arrow[d,equals]\\
\Pi({}_{t}\backslash f^*(H)) \& \Pi(H_{\unsure{t|f}})
\end{tikzcd}}

il suffirait donc de définir

\[
{}_{t}\backslash f^*(H)\longrightarrow H_t,
\]

mais \deleted{\unsure{avec} le diagramme cartésien}

\diagram{\begin{tikzcd}[ampersand replacement=\&,column sep=12mm]
f^*H=H\times_S T\arrow[r,dash]\arrow[d] \& H\arrow[d]\\
S\arrow[r,dash,"f"] \& T\ni t
\end{tikzcd}}

\unsure{induisant|s'en déduit} $f^*H\unsure{\mathrel{-}|\longrightarrow}H$ d'où

\[
\deleted{f^*(H)_t=S_{\unsure{t|f}}\times H_t\xrightarrow{\mathrm{pr}_2}H_t}
\]

\deleted{d'autre part on a}

\[
\deleted{f^*(H)_t\longrightarrow {}_{t}\backslash f^*(H)}
\]

\[
{}_{t}\backslash f^*(H)\longrightarrow {}_{t}\backslash H
\qquad\substack{\text{(au lieu de s'envoyer}\\\text{dans }H_t\text{)}}
\]

mais d'autre part, comme \unsure{$H$|tout} est fibrée sur $T$ donc lisse sur $T$, l'inclusion

\[
\deleted{\unsure{t}}H_t\longrightarrow {}_{t}\backslash H
\]

est une Hot-équivalence. Ou disons plutôt : $f^*H\longrightarrow H$ est un morphisme de catégories sur $T$, induisant

\diagram{\begin{tikzcd}[ampersand replacement=\&,row sep=3mm]
\deleted{\unsure{f_!(f)}}\Pi(f^*H/T)\arrow[r]\arrow[d,equals,"\mathrm{def}"] \& \Pi(H/T),\\
f_!(f^*H) \&
\end{tikzcd}}
